\subsection{The monic pencil polynomial and a cutoff discriminator}
\label{ssec:monic-pencil-cutoff}

The zero-shift pencil supplies a monic extremal problem that is different
from the ground-state Rayleigh quotient and the bordering evaluation
constraint. In the squared coordinate $t=(\nu/\beta)^2$, write
\[
 Q_N(t)=\prod_{j=1}^N(t-j^2),\qquad
 P_a(t)=a_0Q_N(t)+\sum_{j=1}^Nw_jv_j
                \frac{tQ_N(t)}{t-j^2},\qquad w_j=\sqrt2.
\]
This polynomial map from the even coefficient vector $a=(a_0,v)$ is
invertible. Its leading coefficient is $\ell^Ta$, where $\ell=(1,w)$.
Thus degree at most $N-1$ is exactly the endpoint-vanishing space
$a=\mathcal Jx$ used above.

\begin{proposition}[Monic characterization of the zero-shift pencil]
\label{prop:monic-pencil}
Suppose the finite even matrix $A$ is positive definite. In the induced
polynomial inner product $\langle P_a,P_b\rangle=a^TAb$, the degree-$N$
monic orthogonal polynomial has coefficient vector
\[
 a_{\rm OP}=\frac{A^{-1}\ell}{\ell^TA^{-1}\ell}.
\]
Where the central-normalized virtual vector in
Proposition~\ref{prop:arbitrary-shift-pencil} exists at $z=0$, its
polynomial is a nonzero multiple of this monic polynomial.
\end{proposition}
\begin{proof}
The leading-coefficient constraint is $\ell^Ta=1$.
Strict positivity makes its quadratic minimizer unique, with Euler
equation $Aa=\eta\ell$. Every lower-degree polynomial has coefficient
vector in $\ker\ell^T$, proving orthogonality. For the zero-shift
variables, $Bv=-c$, $Bq=w$ and $c^Tq=1-W$ give exactly
\[
 A\binom{1}{v+(\sigma/W)q}=\frac{\sigma}{W}\ell.
\]
After leading-coefficient normalization this is the same minimizer.
\end{proof}
The expression through $A^{-1}\ell$ also avoids dividing by a possibly
vanishing central coordinate. This is a finite algebraic statement;
it does not require RH. It gives the correct monic object for studying
the zero-shift channel, without reinterpreting the ground polynomial
as an ordinary monic minimizer.

\paragraph{The varying measure and its limitation.}
The Gamma-transform identity gives, under a positive real atomic Weil
representation, the explicit measure
\[
 d\mu_{C,N}(t)=\sum_m
 \frac{2L\,m_m\,\delta_{(\gamma_m/\beta)^2}(dt)}
 {\Gamma(N+1+\gamma_m/\beta)^2
  \Gamma(N+1-\gamma_m/\beta)^2},
\]
where $m_m$ is the multiplicity and reciprocal Gamma factors are
interpreted by entire continuation. It represents the polynomial
inner product above. Without the positive real hypothesis, the full
signed-complex Weil functional remains the moment source; finite
matrix positivity alone does not identify it with this positive
zeta-atomic measure. Moments through degree $2N$ converge, which is
all that is needed here. No all-degree finite-moment assumption is made.

In the degree-$(N-1)$ space the pair $H_0,T_0$ is precisely the moment
pair for multiplication by $t$. Its eigenvalues are therefore the
zeros of the degree-$N$ orthogonal polynomial. With distinct ordered
positive-weight atoms $t_k$, min--max gives $\theta_k\ge t_k$;
infinite nondegenerate support gives strictness. Coincident carrier
zeros can remove weights and must not be silently counted. Repeated
ordinates contribute weights rather than independent scalar coordinates.
Thus comparison with a zero list counted with multiplicity needs its
own matching convention.

Factoring $\pi_N(t)=(t-\theta_k)q_k(t)$ and using orthogonality gives
the exact indexed identity
\[
 \theta_k-t_k=
 \frac{\int(t-t_k)q_k(t)^2\,d\mu_{C,N}(t)}
      {\int q_k(t)^2\,d\mu_{C,N}(t)}.
\]
Multiplication by $\beta^2/(\nu_k(0)+\gamma_k)$ converts it to ordinate
error. It exposes the weighted contribution away from the target atom,
but estimating that contribution uniformly is still necessary.
The Gamma weight has the previously derived interior potential
$2N[(1+\tau)\log(1+\tau)+(1-\tau)\log(1-\tau)]$,
$\tau=\gamma\log C/(2\pi N)$. There is now a legitimate monic
object to analyze with that potential. Its varying discrete measure
cannot be replaced by a fixed-measure or Gaussian model without proof.
Rational Ritz convergence through constrained logarithmic potentials
provides relevant methods~\cite{BGV2010}; its asymptotic hypotheses
and individual-index implications remain to be checked here.
The finite polynomial, moment-pencil and virtual-vector identities
pass 283 exact rational checks on positive atomic fixtures.

\paragraph{A new cutoff at the same mode count.}
An independently assembled $C=125,N=1500$ case uses 2300 decimal
assembly digits and 2000 point-solve digits. A 3000-digit prefix control
lies inside all 84 corresponding generator intervals. All first 1000
reference windows are requested before the results are seen, using
the unchanged nearest-three-root predictor and its $1\%$ gate.
Interval roots of the serialized point state and an exact polynomial
upper count give a complete matched prefix of 1000, with no carriers
below its counting threshold. Every prediction passes.
\begin{center}\small
\begin{tabular}{r|rrr}
$C$ ($N=1500$)&$D_1$&$D_{1000}$&Worst relative prediction error\\
\hline
125&662.16446&21.94445&$2.49960\cdot10^{-6}$\\
250&1175.39874&305.89872&$1.69115\cdot10^{-7}$
\end{tabular}
\end{center}
The $C=125$ thousandth zero has more Fourier-band headroom
($\tau\simeq0.72717$, against $0.83156$ at $C=250$), yet far fewer
accurate digits. Mode headroom alone therefore does not determine
accuracy. The cutoff and indexed mode budgets must be distinguished.
The new state has an enclosed matrix residual, but this experiment
does not certify every tiny displacement of the analytic ground
transform. Neither the fixed-parameter successes nor the finite monic
identity close the uniform $C,N,P,k$ law.
